\chapter{Learning with endogeneity: inverse problems and proxy bridges}
\label{r3:inverse-learning}
\section{Which economic quantity must be recovered?}
The invitation experiment may identify a linear coefficient or a complier
mean. A bank asking how funding contribution changes over a continuum of
deposit prices poses a different question. Prices can respond to latent
liquidity demand, so a flexible regression of contribution on price learns
the historical conditional mean, not automatically the structural response.
Equation~\eqref{r3:npiv-operator} instead uses an exclusion and conditional
moment to turn that response into an inverse problem. The same mathematical
issue reappears when historical actions use an unmeasured condition for
which the bank has imperfect proxies.

The analysis proceeds in three steps. First it asks whether the requested
average is identified even when the entire price-response curve is not.
It then shows how regularization trades unstable information against bias.
Finally it compares instruments with proxy measurements, whose identifying
restrictions differ even though both lead to conditional-moment equations.

Suppress observed covariates temporarily. Let
$\mathcal H_A=L_2(P_A)$, $\mathcal H_Z=L_2(P_Z)$ and
$T:\mathcal H_A\to\mathcal H_Z$ be
$(Th)(z)=\E[h(A)\mid Z=z]$. Conditional expectation is a bounded
linear contraction. Its adjoint is
$(T^*q)(a)=\E[q(Z)\mid A=a]$, since
\[
 \langle Th,q\rangle_{P_Z}
 =\E[h(A)q(Z)]=\langle h,T^*q\rangle_{P_A}.
\]
The observed regression $m(z)=\E[Y\mid Z=z]$ obeys $Th=m$ under
the structural conditional moment. Existence says that $m$ is in the
range of $T$ for the chosen model class. Identification of the entire
response says that the kernel $\ker T$ contains only zero on the
allowed difference class. These are separate restrictions. Fitting
an approximate solution is evidence for neither without further analysis.

\subsection{A target can be identified when the whole response is not}
\label{r3:inverse-functional}
Suppose one solution $h_0$ exists. In a linear unrestricted difference
space all solutions are $h_0+v$ for $v\in\ker T$. A linear target
$\theta=\langle\ell,h\rangle_{P_A}$ is therefore identified precisely
when
\begin{equation}
 \langle\ell,v\rangle_{P_A}=0\quad\text{for every }v\in\ker T.
 \label{r3:null-functional}
\end{equation}
This follows by subtracting the targets of two solutions; it is a
project linear-algebra derivation at a fixed observed law. If $h$ is
restricted by shape, bounds or a nonlinear model, the allowed differences
must be restricted as well. The same statement with the full linear
kernel need not be necessary in that smaller class.

For a finite illustration let $A$ take three values and let binary $Z$
supply the two averaging rows
\[
 T=\begin{pmatrix}1/2&1/2&0\\0&1/2&1/2\end{pmatrix}.
\]
Then $T(1,-1,1)'=0$. The responses
$h=(1,2,3)'$ and $h'=(2,1,4)'$ both yield $m=(1.5,2.5)'$.
The equal-weight target $(h_1+h_2+h_3)/3$ differs, being $2$ and
$7/3$. The target with policy weights $(1/4,1/2,1/4)$ equals $2$
in both models, because those weights are the average of the two
observed operator rows. An invitation mix reproducing that distribution
has an identified average even though the response at every price is
not identified. Extrapolating to an equal price mix requires more.

If there exists $q\in\mathcal H_Z$ such that $T^*q=\ell$, then
\begin{equation}
 \theta=\langle\ell,h\rangle
       =\langle q,Th\rangle=\E[q(Z)Y].
 \label{r3:adjoint-target}
\end{equation}
This is an observable representation whose stability depends on $q$.
In infinite dimensions,
$(\ker T)^\perp=\overline{\operatorname{range}(T^*)}$.
Membership in that closure establishes the linear identification
condition, but need not furnish a square-integrable representer in the
range itself. Thus identification of a functional can coexist with a
failure of stable root-$n$ inference. Estimating the operator and
representer adds statistical conditions beyond this Hilbert-space
identity; it cannot be elided by calling an average ``low dimensional.''

\section{Regularization is an assumption about unresolved directions}
\label{r3:inverse-rates}
Assume a compact injective operator with singular system
$T\phi_j=\sigma_j\chi_j$, $\sigma_j>0$. Then
$h=\sum_j h_j\phi_j$ and $m=\sum_j\sigma_jh_j\chi_j$.
The map from a data error in coefficient $j$ to a response error divides
by $\sigma_j$. The instability concerns directions in function space,
not just one scalar first-stage correlation.

In a deliberately simplified Gaussian sequence experiment, suppose
$\widehat m_j=\sigma_jh_j+n^{-1/2}\xi_j$ with independent
standard normal $\xi_j$ and known singular values. A cutoff estimate
retaining $j\le J$ has exact squared-error risk
\begin{equation}
 \E\|\widehat h_J-h\|_2^2
 =\frac1n\sum_{j\le J}\sigma_j^{-2}+\sum_{j>J}h_j^2.
 \label{r3:inverse-risk}
\end{equation}
For $\sigma_j\asymp j^{-a}$ and an ellipsoid
$\sum_j j^{2s}h_j^2\le R^2$, the terms are bounded in order by
$n^{-1}J^{2a+1}$ and $R^2J^{-2s}$. Balancing gives
$J\asymp n^{1/(2s+2a+1)}$ and squared risk of order
$n^{-2s/(2s+2a+1)}$. If $\sigma_j\asymp e^{-cj^a}$, variance
grows exponentially with the cutoff and the corresponding balance is
logarithmic in $n$, giving the slower order $(\log n)^{-2s/a}$.
These calculations expose the tradeoff; they are not an NPIV theorem
for an estimated first stage or a data-selected neural architecture.

For the linear target with coefficients $\ell_j$, the candidate
representer has coefficients $\ell_j/\sigma_j$. In the same sequence
experiment its variance is
$n^{-1}\sum_j\ell_j^2/\sigma_j^2$, when finite. A sufficiently
smooth average can be estimable much faster than the full response.
A sharp threshold or point evaluation can put large weight on exactly
the unstable directions. Consequently the relevant inverse problem
must be evaluated for the bank's requested price-policy contrast,
not inferred from average prediction error.

\subsection{What a primary NPIV rate theorem actually assumes}
\citet[Section 2.1]{hall2005} works with a positive integral operator
formed from the joint density of an endogenous regressor and instrument.
Their operator is a density-weighted normal operator; it is not literally
the conditional-expectation operator $T$ used above. Their eigenvalue
exponent $\alpha$ concerns that positive operator, and must not be
silently read as the singular-value exponent $a$.

In their bivariate kernel result the observations are independent and
identically distributed; the regressor and instrument are supported
on the unit square; their joint density has enough bounded derivatives;
and the relevant conditional second moments are bounded. In the
operator eigenbasis, response coefficients satisfy $|b_j|\le Cj^{-\beta}$,
eigenvalues obey the lower decay bound $j^{-\alpha}\le C\lambda_j$,
and joint-density coefficients obey the summability restriction in A.3.
The exponents satisfy $\alpha>1$, $\beta>1/2$ and
$\beta-1/2\le\alpha<2\beta$. With the density smoothness and
bandwidth restrictions in A.2, A.4 and A.5 and ridge
$a_n\asymp n^{-\alpha/(2\beta+\alpha)}$, their Theorem 4.1 gives
\begin{equation}
 \sup_{G\in\mathcal G}\E_G\|\widehat h-h\|_{L_2([0,1],dx)}^2
       =O\!\left(n^{-(2\beta-1)/(2\beta+\alpha)}\right).
 \label{r3:hall-rate}
\end{equation}
Here the source's norm integrates with respect to Lebesgue measure on
the unit interval; it is not silently the earlier $L_2(P_A)$ norm.
The theorem handles an estimated operator; the sequence illustration
does not. The coefficient-decay condition gives a tail of order
$J^{-(2\beta-1)}$, which explains why its smoothness notation differs
from the ellipsoid exponent $s$. The separate series result imposes
a strong diagonality condition. None of these statements gives a
confidence band merely from choosing a ridge parameter, and the
conditions cannot be transferred to a bank panel with common shocks
without a new sampling argument.

\section{Proxy identification uses different exclusions from IV}
\label{r3:proxy-identification}
Suppose $U$ is latent liquidity need, $X$ are observed pretreatment
covariates, $A$ is the bank's action, and $Z,W$ are two sets of proxies.
Here $Z$ is not the randomized invitation of the earlier example.
Its reused letter denotes a treatment-inducing proxy; $W$ is an
outcome-inducing proxy. A recorded private risk assessment may help
predict $U$ while failing to be an instrument precisely because $U$
also affects $Y$. Proxy identification can permit that dependence,
but it requires different exclusion restrictions.

For each action $a$, assume consistency, positivity given $(U,X)$,
$Y(a)\perp A\mid U,X$, and the proxy restrictions
\[
 Y\perp Z\mid A,U,X,\qquad W\perp(A,Z)\mid U,X.
\]
They can follow from a suitable latent DAG with no effect of $A,Z$
on the pretreatment $W$ and no direct effect of $Z$ on $Y$.
Using a post-assistance balance as $W$ would ordinarily violate this
restriction. The latent condition $U$ must include the relevant shared
causes; labelling two arbitrary database columns as proxies does not
establish it.

Suppose a square-integrable outcome bridge exists satisfying
\begin{equation}
 \E[Y\mid Z,A=a,X]=\E[h(W,a,X)\mid Z,A=a,X].
 \label{r3:outcome-bridge}
\end{equation}
Also assume latent completeness: for fixed $a,x$,
$\E[v(U)\mid Z,A=a,X=x]=0$ implies $v(U)=0$ almost surely
in the relevant class. Define
$d_a(U,X)=\E[Y\mid U,A=a,X]-\E[h(W,a,X)\mid U,X]$.
The exclusions and bridge equation give
$\E[d_a(U,X)\mid Z,A=a,X]=0$. Completeness makes $d_a=0$.
Consistency and latent exchangeability then give
\begin{align}
 \E[Y(a)]&=\E[\E[Y\mid U,A=a,X]]\nonumber\\
 &=\E[\E[h(W,a,X)\mid U,X]]
  =\E[h(W,a,X)].
 \label{r3:proximal-mean}
\end{align}
This is a mean-level reconstruction of the proxy-identification
argument in \citet[Sections 2--3 and Theorem 1]{miao2018}; observed
covariates and the mean-bridge formulation are stated explicitly in
\citet[Section 4, Assumptions 4--6 and Theorem 3]{ghassami2022}.

Three restrictions must remain distinct. Latent completeness justifies
the step from an observable bridge equation to a latent causal equality.
Bridge existence says that the observable conditional mean lies in an
operator range. Observable completeness of $W$ given $(Z,A,X)$ can
give uniqueness of $h$. A matrix rank estimated from proxies does not
directly verify latent completeness; moreover uniqueness of every
bridge value is stronger than uniqueness of its causal average.
\citet{miao2018} explicitly permits nonunique bridge solutions while
identifying the effect under the stated restrictions.

\subsection{A finite banking construction with an exact answer}
Let $P(U=1)=1/2$. Conditional on $U$, generate independent $A,Z,W$
with probabilities of one
\[
 p_A(U)=1/5+3U/5,\quad
 p_Z(U)=3/10+2U/5,\quad
 p_W(U)=1/10+4U/5,
\]
and let $Y=2A+3U$. All combinations have positive probability.
The causal effect is $2$. The raw treated-minus-untreated contrast
is $2+3(0.8-0.2)=3.8$. Adjusting for a single noisy proxy need not
recover $2$.

The outcome bridge
\begin{equation}
 h(w,a)=2a-3/8+15w/4
 \label{r3:proxy-example-bridge}
\end{equation}
satisfies $\E[h(W,a)\mid U]=2a+3U$, hence the observable
bridge equation. Since $P(W=1)=1/2$, averaging gives
$\E[h(W,a)]=2a+1.5$ and effect $2$.
The $2\times2$ matrix with rows $P(W=\cdot\mid Z=z,A=a)$
is invertible in either action stratum, so this bridge can be recovered
from observed conditional moments without observing $U$.

Now let $W$ lose all information about $U$ while retaining independent
noise. Its conditional rows become equal. Full column rank is lost and the
bridge cannot reproduce a $Z$-varying outcome regression. Replacing
the inverse by a pseudoinverse returns a numerical vector, but does not
restore the missing equation or identify the original target. This is
the proxy analogue of a weak or unidentified IV direction.

\section{Two bridges produce a doubly robust estimating expression}
\label{r3:proximal-score}
Suppose a treatment bridge also exists with
\begin{equation}
 \E[q(Z,a,X)\mid W,A=a,X]
       =\frac1{P(A=a\mid W,X)}.
 \label{r3:treatment-bridge}
\end{equation}
Together with the latent restrictions needed for proximal identification,
consider
\begin{equation}
 \Phi_a(h,q)=h(W,a,X)+\ind\{A=a\}q(Z,a,X)
                                  \{Y-h(W,a,X)\}.
 \label{r3:proximal-augmented}
\end{equation}
If $h=h_0$ is a valid outcome bridge, the residual term has mean zero
for any integrable $q$ measurable in $(Z,a,X)$, by conditioning on
$(Z,A,X)$. If $q=q_0$ is a valid treatment bridge, then for any $h$
\[
 \E[\ind\{A=a\}q_0(Z,a,X)h(W,a,X)]=\E[h(W,a,X)].
\]
The $h$ terms cancel, and the expectation equals
$\E[\ind\{A=a\}q_0Y]=\E[Y(a)]$ when the compatible true
outcome bridge exists. This proves the union-model robustness of
the estimating expression at fixed candidates.

Write $\delta_h=\widetilde h-h_0$ and
$\delta_q=\widetilde q-q_0$. Expanding and using both bridge
equations gives the exact remainder
\begin{equation}
 \E[\Phi_a(\widetilde h,\widetilde q)]-\E[Y(a)]
 =-\E[\ind\{A=a\}\delta_q(Z,a,X)\delta_h(W,a,X)].
 \label{r3:proximal-remainder}
\end{equation}
Conditioning this product on $(Z,A=a,X)$ controls it by a norm of
$\delta_q$ times a norm of the projected bridge error
$\E[\delta_h\mid Z,A=a,X]$. Conditioning in the other direction
gives the complementary bound. An inverse problem can have small
projected error but large unprojected error. Both directions therefore
have to be evaluated using the operator's ill-posedness, rather than
copying the AIPW nuisance-rate rule from an ordinary regression problem.
\citet[Section 3.2, Assumptions 2--3 and Theorem 2]{ghassami2022}
states the relevant consistency, projected-error and ill-posedness
conditions for its cross-fitted asymptotically linear estimator.
Its Section 4 identifies the locally efficient influence function
under the intersection model and uniqueness conditions; those
qualifications are stronger than the algebraic mean identity alone.

\section{Learning the bridge is a conditional-moment problem}
\label{r3:bridge-learning}
For one action stratum let
$r_h(Z,X)=\E[Y-h(W,a,X)\mid Z,A=a,X]$. In the unrestricted
square-integrable test-function class,
\begin{align}
 \sup_f\ \E[\ind\{A=a\}
       \{2f(Z,X)(Y-h(W,a,X))-f(Z,X)^2\}]
 =\E[\ind\{A=a\}r_h(Z,X)^2].
 \label{r3:minimax-moment}
\end{align}
Complete the square conditionally to prove the identity: the maximizer
is $f=r_h$. Minimizing the left side over $h$ targets the conditional
moment rather than individual response prediction. Restricting $f$ to
an RKHS or neural class changes which moment violations are detected;
regularizing $h$ selects among unstable approximate solutions.
\citet[Sections 3 and 5]{ghassami2022} develops RKHS minimax
constructions and rate conditions for this class of problems.

In contrast, least-squares regression of $Y$ on $(W,A,X)$ learns
$\E[Y\mid W,A,X]$, which is generally not an outcome bridge. In
the exact example, posterior $U$ given $W,A$ differs from the linear
bridge correction in Equation~\eqref{r3:proxy-example-bridge}. A
learner can therefore have excellent factual prediction and fail the
observable bridge moments. Validation needs held-out moment violations,
regularization sensitivity and target uncertainty in addition to factual
loss. Even zero observed violations cannot verify the unobserved
exclusions; they only test implications within the chosen function class.

For large banking data, a finite-basis conditional-moment model leads to
cross-products of instrument and response bases, an ill-conditioned
linear system, and a regularized solve. Dense kernel methods instead
form matrices whose memory is quadratic in sample size. Approximations
using finite features or blocks change the numerical moment problem;
their residual error must be propagated into the target remainder.
An architecture decision follows from these mathematical objects,
not from declaring every inverse learner interchangeable with a
supervised regression adapter.

The proposed product should return the target functional, operator and
proxy roles, rank/completeness assumptions, fitted moment residuals,
regularization path, and the exact statistical conditions claimed.
It should distinguish a nonunique structural response with an identified
average from an unidentified average, and distinguish both from an
identified but statistically unstable target. These distinctions also
prepare the partially observed longitudinal problem: a hidden-state
likelihood, a proxy bridge and a causal value solve different recovery
problems even when they share the same observations.
